\section*{Reproducibility, implementation and data availability}
\label{sec:reproducibility}

\paragraph{Code and data availability.}
The complete source code and post-processing scripts used in this work are hosted on GitHub at \url{https://github.com/USERNAME/REPOSITORY} and permanently archived on Zenodo at \url{https://doi.org/10.5281/zenodo.XXXXXXX}.
The Zenodo record corresponds to the exact version used for this manuscript (release tag and commit hash recorded in the archive metadata) and includes (i) all JSON parameter files, (ii) the anatomy- and gait-related inputs required for the femur case study, and (iii) the raw simulation outputs from which the figures were generated.
The repository is organized so that every simulation run is self-describing: it writes a complete \texttt{config.json} snapshot, rank-0 logs, and machine-readable CSV metrics (\texttt{steps.csv}, \texttt{subiterations.csv}) into the output directory together with the field data.

\paragraph{Software stack and execution.}
The solver is implemented in Python on top of FEniCSx/dolfinx and PETSc, and all PDE blocks are solved in distributed memory via MPI.
Field output is written collectively using the ADIOS2-backed VTX format (\texttt{*.bp}) for scalable parallel I/O and restart/checkpointing.
A minimal end-to-end example is the box benchmark, which can be run as
\texttt{mpirun -n 4 python run\_box\_model.py},
producing a self-contained results directory (including \texttt{config.json} and all logs/outputs).
All parameter sweeps and figure-generation scripts used in the manuscript are provided in the repository (see \texttt{run\_*.py} and \texttt{analysis/*.py}).

\paragraph{Modular staggered architecture and extensibility.}
The coupled problem is advanced by a partitioned block Gauss--Seidel scheme (mechanics $\rightarrow$ fabric $\rightarrow$ stimulus $\rightarrow$ density; \S\ref{sec:fp-aa}).
In the implementation, each subproblem is encapsulated as an independent \emph{coupling block} exposing its state and output fields (see \texttt{simulation/protocols.py} and \texttt{simulation/registry.py}).
Blocks are registered in \texttt{simulation/model.py}; the fixed-point solver, time integrator, I/O, and diagnostics operate on the registry and therefore remain agnostic to the specific set of fields.
Adding a new field to the staggered scheme requires only (i) defining the new \texttt{dolfinx.fem.Function} and its ``old'' counterpart, (ii) implementing one additional block that advances it, and (iii) registering the block in the desired Gauss--Seidel order; the rest of the workflow (time stepping, logging, output) is discovered automatically from the registry.

\paragraph{Parallel execution and scaling.}
All simulations are MPI-parallel and require no shared-memory assumptions.
For parameter studies we additionally support concurrent execution of multiple independent runs by splitting \texttt{MPI\_COMM\_WORLD} into subcommunicators (see \texttt{parametrizer.py}), enabling both strong scaling (single large run) and throughput scaling (many runs in parallel).
To document parallel performance, we include a strong scaling benchmark on the stiff femur case study (see \texttt{run\_weak\_scaling.py}).
Executing the fully coupled solver on this representative mesh (approx.~75k elements, 40k mechanics DOFs) demonstrates super-linear scaling up to 2 cores and maintains 56\% efficiency at 8 cores (speedup $4.45\times$), confirming that the partitioned solver effectively utilizes workstation-level parallelism even for moderate problem sizes (\cref{fig:strong_scaling}).

\begin{figure}[htbp]
    \centering
    \safeincludegraphics[width=0.6\linewidth]{strong_scaling.png}
    \caption{Strong scaling of the coupled bone remodeling solver on the stiff femur configuration ($N_{\mathrm{cells}}\approx 75\mathrm{k}$, $N_{\mathrm{dofs}}^{\mathrm{mech}}\approx 40\mathrm{k}$). The solver achieves super-linear speedup at 2 cores (likely due to cache effects) and scales with reasonable efficiency up to 8 cores despite the relatively small problem size per core.}
    \label{fig:strong_scaling}
\end{figure}

\paragraph{Verification.}
Core infrastructure (time integration, MPI-consistent diagnostics, I/O) is covered by automated tests in \texttt{tests/} runnable via \texttt{pytest} (serial or under MPI).
